\documentclass[11pt,a4paper]{article}
\usepackage[utf8]{inputenc}
\usepackage{amsmath,amssymb,amsthm}
\usepackage{listings}
\usepackage[margin=2.5cm]{geometry}
\usepackage{xcolor}

\lstset{
  language=C++,
  basicstyle=\ttfamily\small,
  keywordstyle=\color{blue}\bfseries,
  commentstyle=\color{green!50!black},
  stringstyle=\color{red!70!black},
  numbers=left,
  numberstyle=\tiny\color{gray},
  breaklines=true,
  frame=single,
  tabsize=4,
  showstringspaces=false
}

\title{IOI 1992 -- Problem 3: Islands}
\author{Solution}
\date{}

\begin{document}
\maketitle

\section{Problem Statement}

Given an $R \times C$ elevation grid and a water level $L$, a cell is
\emph{land} if its elevation is strictly greater than $L$, and \emph{water}
otherwise. Count the number of islands (maximal connected components of land
cells under 4-adjacency).

An extended version asks for island counts at multiple water levels.

\section{Solution}

\subsection{Single Water Level: Flood Fill}

\begin{enumerate}
  \item Mark each cell as land if $\text{elevation}[r][c] > L$.
  \item Count connected components of land cells using BFS.
\end{enumerate}

\subsection{Multiple Water Levels: Offline DSU}

To answer $Q$ queries efficiently:
\begin{enumerate}
  \item Sort all cells by elevation in decreasing order.
  \item Process cells one by one. When a cell ``emerges'' above the current
    water level, activate it in a Disjoint Set Union (DSU) structure and
    merge it with any already-active 4-neighbors.
  \item Record the component count at each distinct elevation threshold.
  \item Answer each query by looking up the component count at the
    appropriate threshold.
\end{enumerate}

\section{Complexity Analysis}

\begin{itemize}
  \item \textbf{Single water level:} $O(RC)$ time and space.
  \item \textbf{Multiple water levels:} $O(RC\,\alpha(RC) + Q\log(RC))$
    total, where $\alpha$ is the inverse Ackermann function.
\end{itemize}

\section{Implementation: Single Water Level}

\begin{lstlisting}
#include <iostream>
#include <vector>
#include <queue>
using namespace std;

int main() {
    int R, C, L;
    cin >> R >> C >> L;

    vector<vector<int>> elev(R, vector<int>(C));
    for (int i = 0; i < R; i++)
        for (int j = 0; j < C; j++)
            cin >> elev[i][j];

    vector<vector<bool>> visited(R, vector<bool>(C, false));
    const int dx[] = {0, 0, 1, -1};
    const int dy[] = {1, -1, 0, 0};
    int islands = 0;

    for (int i = 0; i < R; i++) {
        for (int j = 0; j < C; j++) {
            if (elev[i][j] > L && !visited[i][j]) {
                islands++;
                queue<pair<int,int>> q;
                q.push({i, j});
                visited[i][j] = true;
                while (!q.empty()) {
                    auto [r, c] = q.front(); q.pop();
                    for (int d = 0; d < 4; d++) {
                        int nr = r + dx[d], nc = c + dy[d];
                        if (nr >= 0 && nr < R && nc >= 0 && nc < C
                            && !visited[nr][nc]
                            && elev[nr][nc] > L) {
                            visited[nr][nc] = true;
                            q.push({nr, nc});
                        }
                    }
                }
            }
        }
    }

    cout << islands << endl;
    return 0;
}
\end{lstlisting}

\section{Implementation: Multiple Water Levels (DSU)}

\begin{lstlisting}
#include <iostream>
#include <vector>
#include <algorithm>
#include <map>
using namespace std;

struct DSU {
    vector<int> parent, rank_;
    int components;

    DSU(int n) : parent(n, -1), rank_(n, 0), components(0) {}

    void activate(int x) {
        parent[x] = x;
        components++;
    }

    bool active(int x) const { return parent[x] != -1; }

    int find(int x) {
        while (parent[x] != x) {
            parent[x] = parent[parent[x]]; // path splitting
            x = parent[x];
        }
        return x;
    }

    void unite(int a, int b) {
        a = find(a); b = find(b);
        if (a == b) return;
        if (rank_[a] < rank_[b]) swap(a, b);
        parent[b] = a;
        if (rank_[a] == rank_[b]) rank_[a]++;
        components--;
    }
};

int main() {
    int R, C;
    cin >> R >> C;

    vector<vector<int>> elev(R, vector<int>(C));
    vector<pair<int,int>> cells; // (elevation, flat_index)

    for (int i = 0; i < R; i++)
        for (int j = 0; j < C; j++) {
            cin >> elev[i][j];
            cells.push_back({elev[i][j], i * C + j});
        }

    // Sort by decreasing elevation
    sort(cells.begin(), cells.end(), greater<>());

    DSU dsu(R * C);
    const int dx[] = {0, 0, 1, -1};
    const int dy[] = {1, -1, 0, 0};

    // Process cells from highest to lowest elevation.
    // After processing all cells with elevation > L,
    // dsu.components gives the island count at water level L.
    map<int, int> levelIslands;

    int i = 0;
    while (i < (int)cells.size()) {
        int e = cells[i].first;
        // Process all cells with this elevation
        int j = i;
        while (j < (int)cells.size() && cells[j].first == e) {
            int idx = cells[j].second;
            int r = idx / C, c = idx % C;
            dsu.activate(idx);
            for (int d = 0; d < 4; d++) {
                int nr = r + dx[d], nc = c + dy[d];
                if (nr >= 0 && nr < R && nc >= 0 && nc < C) {
                    int nidx = nr * C + nc;
                    if (dsu.active(nidx))
                        dsu.unite(idx, nidx);
                }
            }
            j++;
        }
        // Water level = e - 1 means all cells with elev > e-1
        // (i.e., elev >= e) are land.
        levelIslands[e - 1] = dsu.components;
        i = j;
    }

    int Q;
    cin >> Q;
    while (Q--) {
        int L;
        cin >> L;
        auto it = levelIslands.upper_bound(L);
        if (it == levelIslands.begin())
            cout << 0 << "\n";
        else {
            --it;
            cout << it->second << "\n";
        }
    }

    return 0;
}
\end{lstlisting}

\section{Notes}

\begin{itemize}
  \item The single-water-level solution is a straightforward BFS flood fill.
  \item The DSU solution processes cells in decreasing elevation order. When
    a batch of cells at elevation~$e$ is activated, the resulting component
    count equals the number of islands for any water level in $[e-1, e')$
    where $e'$ is the next lower elevation present.
  \item Path splitting (iterative path compression) is used instead of
    recursive path compression to avoid stack overflow on large inputs.
\end{itemize}

\end{document}
